\documentclass[11pt,a4paper]{article}
\usepackage[utf8]{inputenc}
\usepackage[T1]{fontenc}
\usepackage{lmodern}
\usepackage[french]{babel}
\usepackage{amsmath,amssymb,mathtools}
\usepackage{icomma}
\usepackage{xcolor}
\usepackage{tikz}
\usepackage{pgfplots}
\usepackage[most]{tcolorbox}
\usepackage{enumitem}
\usepackage{fancyhdr}
\usepackage[hidelinks]{hyperref}
\usepackage{titlesec}
\usepackage[a4paper,margin=2.2cm,headheight=26pt,headsep=10pt]{geometry}
\usetikzlibrary{arrows.meta,positioning,calc,intersections,angles,quotes}
\usepgfplotslibrary{fillbetween}
\pgfplotsset{compat=1.18}
\definecolor{primary}{RGB}{30,75,145}
\definecolor{primarydark}{RGB}{20,50,100}
\definecolor{defcol}{RGB}{0,114,178}
\definecolor{thmcol}{RGB}{0,135,110}
\definecolor{excol}{RGB}{0,130,85}
\definecolor{remarkcol}{RGB}{120,70,180}
\definecolor{attncol}{RGB}{200,40,55}
\definecolor{methcol}{RGB}{85,85,100}
\definecolor{areacol}{RGB}{120,160,230}
\definecolor{areacol2}{RGB}{230,150,80}
\newcommand{\dd}{\mathrm{d}}
\newcommand{\R}{\mathbb{R}}
\newcommand{\ua}{\,\text{u.a.}}
\newcommand{\tg}{\tan}\newcommand{\cotg}{\cot}
\tcbset{boxbase/.style={enhanced,breakable,boxrule=0.9pt,arc=2.5pt,left=3mm,right=3mm,top=2.3mm,bottom=2.3mm,fonttitle=\bfseries,coltitle=white,toptitle=0.6mm,bottomtitle=0.6mm}}
\newtcolorbox{methode}[1][]{boxbase,colback=methcol!7,colframe=methcol,sharp corners,title={\textsc{Comment faire}\ifx\relax#1\relax\relax\else\textmd{ : \itshape #1}\fi}}
\newtcolorbox{exemple}[1][]{boxbase,colback=excol!5,colframe=excol,title={\textsc{Exemple}\ifx\relax#1\relax\relax\else\textmd{ : \itshape #1}\fi}}
\newtcolorbox{idee}[1][]{boxbase,colback=defcol!5,colframe=defcol,title={\textsc{Idée clé}\ifx\relax#1\relax\relax\else\textmd{ : \itshape #1}\fi}}
\newtcolorbox{remarque}[1][]{boxbase,colback=remarkcol!6,colframe=remarkcol,title={\textsc{Remarque}\ifx\relax#1\relax\relax\else\textmd{ : \itshape #1}\fi}}
\newtcolorbox{attention}[1][]{boxbase,colback=attncol!6,colframe=attncol,title={\textsc{Attention}\ifx\relax#1\relax\relax\else\textmd{ : \itshape #1}\fi}}
\newtcolorbox{exo}[1][]{boxbase,colback={primary!4},colframe=primary,title={\textsc{Exercice #1}}}
\newtcolorbox{corr}[1][]{boxbase,colback={thmcol!5},colframe=thmcol,title={\textsc{Corrigé #1}}}
\tikzset{axe/.style={->,>=Stealth,black!65,line width=0.7pt},courbe/.style={primary,line width=1.3pt},courbe2/.style={areacol2!90!black,line width=1.3pt},dvert/.style={gray,dashed,line width=0.7pt}}
\pagestyle{fancy}\fancyhf{}
\fancyhead[L]{\small\textcolor{primarydark}{\textbf{Fiche 7} — Intégrale d'une fonction continue}}
\fancyhead[R]{\small\textcolor{primarydark}{\textsc{Thème 3 — Fractions rationnelles}}}
\fancyfoot[C]{\small\thepage}
\renewcommand{\headrulewidth}{0.8pt}
\titleformat{\section}{\large\bfseries\color{primarydark}}{\thesection}{0.6em}{}
\begin{document}
\section*{Thème 3 — Fractions rationnelles \quad\textbullet\quad Niveau Difficile}

\begin{center}\itshape
Intégrales définies : décomposer, primitiver, évaluer aux bornes, puis regrouper les $\ln$.\\
On vérifiera à chaque fois que les pôles ne sont pas dans l'intervalle d'intégration.
\end{center}

\begin{exo}[1 — guidé]
On veut calculer $\displaystyle I=\int_{2}^{3} \frac{2x}{(x+1)(x+3)}\,\dd x$.
\begin{enumerate}[label=\alph*),leftmargin=1.8em,topsep=2pt,itemsep=2pt]
  \item Décomposer $\dfrac{2x}{(x+1)(x+3)}$ en fractions simples.
  \item En déduire une primitive $F$ sur $[2,3]$.
  \item Calculer $I=F(3)-F(2)$ et l'écrire sous la forme d'un unique logarithme $\ln\!\big(\tfrac{p}{q}\big)$.
\end{enumerate}
\end{exo}

\begin{exo}[2]
Calculer $\displaystyle \int_{3}^{4} \frac{1}{(x-1)(x-2)}\,\dd x$ et donner le résultat sous la
forme d'un seul logarithme.
\end{exo}

\begin{exo}[3]
Calculer $\displaystyle \int_{0}^{1} \frac{4}{(x+1)(x+3)}\,\dd x$ et regrouper le résultat en un
seul logarithme.
\end{exo}

\begin{exo}[4 — guidé, avec factorisation]
On pose $\displaystyle J=\int_{1}^{2} \frac{x+3}{x^2+3x+2}\,\dd x$.
\begin{enumerate}[label=\alph*),leftmargin=1.8em,topsep=2pt,itemsep=2pt]
  \item Factoriser $x^2+3x+2$ et décomposer $\dfrac{x+3}{x^2+3x+2}$ en fractions simples.
  \item Déterminer une primitive $G$ sur $[1,2]$.
  \item Calculer $J=G(2)-G(1)$ et l'exprimer sous la forme $\ln\!\big(\tfrac{p}{q}\big)$.
\end{enumerate}
\end{exo}

\end{document}
